%---------------------------------------------------------------------
% Urdu edition: formula sheet and the recurring mistakes.
%---------------------------------------------------------------------

\section*{فارمولا شیٹ}
\markboth{فارمولا شیٹ}{}

نو یونٹوں کا سارا مواد دو صفحوں پر۔ ہر اصول کی \emph{شکل} اور اس کے استعمال کا
موقع یاد رکھیے؛ حساب آسان حصہ ہے۔

\subsection*{احتمال (یونٹ \textenglish{1--2})}

\begin{center}
\begin{tabular}{@{}p{4.0cm} p{5.6cm} p{4.6cm}@{}}
\toprule
\textbf{اصول} & \textbf{فارمولا} & \textbf{کب استعمال کریں} \\
\midrule
متمم        & \(\prob{A'}=1-\prob{A}\) & ''کم از کم ایک''، ''نہیں'' \\
جمع (عام)   & \(\prob{A\cup B}=\prob{A}+\prob{B}-\prob{A\cap B}\) & ''یا''، واقعات میں اشتراک \\
جمع (خارج)  & \(\prob{A\cup B}=\prob{A}+\prob{B}\) & ''یا''، \(A\cap B=\varnothing\) \\
ضرب (آزاد)  & \(\prob{A\cap B}=\prob{A}\prob{B}\) & ''اور''، واپسی کے ساتھ \\
ضرب (وابستہ)& \(\prob{A\cap B}=\prob{A}\prob{B\mid A}\) & ''اور''، بغیر واپسی \\
مشروط       & \(\prob{A\mid B}=\dfrac{\prob{A\cap B}}{\prob{B}}\) & ''بشرطیکہ'' \\
آزادی کی جانچ & \(\prob{A\cap B}=\prob{A}\prob{B}\) & آزادی ثابت یا رد کرنے کو \\
\midrule
دو ذی حدی   & \(P(X=x)=\dbinom{n}{x}p^{x}(1-p)^{n-x},\ \ E(X)=np\) & ثابت \(p\)، واپسی کے ساتھ \\
\textenglish{Hypergeometric} & \(P(X=x)=\dfrac{\binom{K}{x}\binom{N-K}{n-x}}{\binom{N}{n}}\) & محدود تھیلا، بغیر واپسی \\
ہر منقطع تقسیم & \(\sum_x P(X=x)=1,\ \ E(X)=\sum_x x\,P(X=x)\) & ہمیشہ --- پڑتال کے طور پر \\
\bottomrule
\end{tabular}
\end{center}

\subsection*{مساوات اور خطوط (یونٹ \textenglish{3--4})}

\begin{center}
\begin{tabular}{@{}p{4.6cm} p{9.8cm}@{}}
\toprule
دو درجی فارمولا & \(x=\dfrac{-b\pm\sqrt{b^{2}-4ac}}{2a}\)؛ ممیز \(>0\) دو جڑیں، \(=0\) ایک، \(<0\) کوئی حقیقی نہیں \\
مربع مکمل کرنا & \(x^{2}+bx=\left(x+\tfrac b2\right)^{2}-\tfrac{b^{2}}{4}\) \\
عدم مساوات & منفی عدد سے ضرب یا تقسیم پر علامت \emph{الٹ} جاتی ہے \\
مطلق قدر (مساوات) & \(|u|=|v| \iff u=v\) یا \(u=-v\) \\
مطلق قدر (\(\le\)) & \(|u|\le a \iff -a\le u\le a\) --- ایک وقفہ \\
مطلق قدر (\(\ge\)) & \(|u|\ge a \iff u\ge a\) یا \(u\le -a\) --- دو وقفے \\
فاصلہ & \(d=\sqrt{(x_2-x_1)^{2}+(y_2-y_1)^{2}}\) \\
وسطی نقطہ & \(M=\left(\dfrac{x_1+x_2}{2},\ \dfrac{y_1+y_2}{2}\right)\) \\
میلان & \(m=\dfrac{y_2-y_1}{x_2-x_1}\) \\
نقطہ--میلان & \(y-y_1=m(x-x_1)\)؛ \quad \(y=mx+c\) \\
متوازی / عمود & \(m_1=m_2\)؛ \quad \(m_1m_2=-1\) (عمودی خطوط پر لاگو نہیں) \\
خطی لاگت & \(C(x)=(\text{فی اکائی متغیر})x+(\text{ثابت لاگت})\) \\
توازن & طلب \(=\) رسد رکھ کر \(q\) نکالیے، پھر \(p\) میں تعویض \\
\bottomrule
\end{tabular}
\end{center}

\subsection*{میٹرکس اور محدد (یونٹ \textenglish{5--6})}

\begin{center}
\begin{tabular}{@{}p{4.6cm} p{9.8cm}@{}}
\toprule
ابعاد کا میل & \(A_{m\times n}B_{n\times p}\) ممکن ہے اور نتیجہ \(m\times p\) \\
منقلب کے اصول & \((A^{t})^{t}=A,\ \ (A+B)^{t}=A^{t}+B^{t},\ \ \boxed{(AB)^{t}=B^{t}A^{t}}\) \\
معکوس کے اصول & \((A^{-1})^{-1}=A,\ \ \boxed{(AB)^{-1}=B^{-1}A^{-1}}\) \\
\(2\times2\) محدد & \(\dt{a&b\\c&d}=ad-bc\) \\
\(2\times2\) معکوس & \(\dfrac{1}{ad-bc}\mat{d&-b\\-c&a}\) \\
ہم عامل & \(C_{ij}=(-1)^{i+j}M_{ij}\)؛ علامتیں \(\mat{+&-&+\\-&+&-\\+&-&+}\) \\
محدد & \(\det A=\sum_j a_{ij}C_{ij}\) کسی بھی سطر یا ستون سے \\
ملحقہ & \(\operatorname{adj}A=\left[C_{ij}\right]^{t}\) --- ہم عاملوں کا \emph{منقلب} \\
معکوس & \(A^{-1}=\dfrac{1}{\det A}\operatorname{adj}A\)، شرط \(\det A\neq0\) \\
کریمر کا اصول & \(x_k=\dfrac{D_k}{D}\)؛ \(D_k\) میں \(k\)واں ستون \(b\) سے بدلیں \\
\(\det=0\) کا مطلب & منفرد حل نہیں: یا کوئی حل نہیں، یا لامتناہی \\
محدد کا ضربی اصول & \(\det(AB)=\det A\cdot\det B\) \\
\bottomrule
\end{tabular}
\end{center}

\subsection*{کیلکولس (یونٹ \textenglish{7--9})}

\begin{center}
\begin{tabular}{@{}p{4.6cm} p{9.8cm}@{}}
\toprule
قوت & \(\dd{}{x}x^{n}=nx^{n-1}\) --- جڑوں اور معکوسوں کو پہلے قوت میں لکھیے \\
حاصل ضرب & \((uv)'=u'v+uv'\) \\
تقسیم & \(\left(\dfrac uv\right)'=\dfrac{u'v-uv'}{v^{2}}\) --- \emph{اوپر} والے کا مشتق پہلے \\
زنجیری & \(\bigl(f(g(x))\bigr)'=f'(g(x))\,g'(x)\) --- باہر پہلے، پھر اندر سے ضرب \\
تفاعلی & \(\dd{}{x}e^{u}=e^{u}u'\)، \quad \(\dd{}{x}\ln u=\dfrac{u'}{u}\) \\
\(\infty\) پر حد & نسب نما کی سب سے بڑی قوت سے تقسیم کیجیے \\
\(\tfrac00\) کی صورت & تحلیل اور منسوخی، پھر دوبارہ تعویض \\
حد کا وجود & \(\lim_{x\to a^{-}}f=\lim_{x\to a^{+}}f\) \\
جزوی مشتق & ایک متغیر میں تفرق، دوسرا ثابت \\
\textenglish{Clairaut} & \(f_{xy}=f_{yx}\) --- مخلوط مشتقوں کی مفت پڑتال \\
ساکن نقطہ & \(f'(x)=0\) (یا دو متغیر میں \(f_x=f_y=0\)) \\
دوسرے مشتق کی جانچ & \(f''<0\) اعظم، \(f''>0\) اصغر، \(f''=0\) بے نتیجہ \\
دو متغیر کی جانچ & \(D=f_{xx}f_{yy}-(f_{xy})^{2}\)؛ \(D>0\) و \(f_{xx}<0\) اعظم، \(f_{xx}>0\) اصغر؛ \(D<0\) زینی \\
انعطافی نقطہ & \(f''=0\) \emph{اور} \(f''\) کی علامت بدلے \\
آمدنی، منافع & \(R=x\,p(x)\)، \quad \(P=R-C\)؛ اعظم جہاں \(MR=MC\) اور \(P''<0\) \\
اوسط لاگت & \(\bar{C}=\dfrac{C(x)}{x}\)؛ اصغر پر \(MC=\bar{C}\) \\
\bottomrule
\end{tabular}
\end{center}

\clearpage

%---------------------------------------------------------------------
\section*{آٹھ عام غلطیاں}
\markboth{آٹھ عام غلطیاں}{}

یہ فہرست اس کتابچے کے تمام سرخ خانوں کا نچوڑ ہے، یونٹوں کی ترتیب سے:

\begin{enumerate}
  \item \textbf{اشتراک والے واقعات پر \(\prob{A}+\prob{B}\) لگانا۔} سادہ جمع کا
        اصول صرف باہم خارج واقعات پر چلتا ہے؛ ورنہ اشتراک گھٹانا پڑتا ہے۔
  \item \textbf{''صرف ایک'' میں گنتی بھول جانا۔} \(\binom{n}{1}\) یہ گنتا ہے کہ
        وہ ایک \emph{کون} ہے۔ ''سب'' اور ''کوئی نہیں'' میں یہ ضرب نہیں لگتی۔
  \item \textbf{''بغیر واپسی'' پر دو ذی حدی لگا دینا۔} یہ الفاظ
        \textenglish{hypergeometric} کی طرف اشارہ ہیں اور جواب واضح طور پر بدل
        دیتے ہیں۔
  \item \textbf{منفی عدد سے ضرب یا تقسیم پر عدم مساوات کی علامت نہ الٹنا۔}
  \item \textbf{\((AB)^{t}=A^{t}B^{t}\) اور \((AB)^{-1}=A^{-1}B^{-1}\) لکھنا۔}
        دونوں میں ترتیب الٹ جاتی ہے۔ دونوں پورے سوال کے برابر ہیں۔
  \item \textbf{ملحقہ بناتے وقت ہم عاملوں کا منقلب بھول جانا۔}
        \(AA^{-1}=I\) کی پڑتال کسی \emph{غیر قطری} خانے پر کیجیے، صرف قطری پر
        نہیں۔
  \item \textbf{کریمر لگانے سے پہلے \(D\) نہ نکالنا۔} \(D=0\) ہو تو اصول لاگو
        نہیں ہوتا اور درست جواب نظام کے بارے میں ایک بیان ہوتا ہے۔
  \item \textbf{\(f'=0\) پر رک جانا۔} دوسرے مشتق کی جانچ ہی اعظم، اصغر اور زینی
        نقطے میں فرق کرتی ہے، اور اس کے نمبر الگ ملتے ہیں۔
\end{enumerate}

%---------------------------------------------------------------------
\section*{اگر وقت کم ہو}

یونٹ \textenglish{4}، \textenglish{5}، \textenglish{6}، \textenglish{7} اور
\textenglish{9} مل کر زیادہ تر نمبر اور تمام طویل سوالات اٹھاتے ہیں۔ ان میں سب
سے زیادہ نفع بخش پانچ مہارتیں یہ ہیں:

\begin{enumerate}
  \item \(3\times3\) نظام پر کریمر کا اصول، بشمول \(D=0\) والی صورت۔
  \item ہم عاملوں سے \(3\times3\) کا معکوس، ملحقہ کے منقلب سمیت۔
  \item حاصل ضرب، تقسیم اور زنجیری اصول --- خاص طور پر تقسیم کے اندر زنجیری۔
  \item منافع، آمدنی یا اوسط لاگت کی بہتر بندی، دوسرے مشتق کی جانچ کے ساتھ۔
  \item دو نقطوں سے، یا ایک نقطے اور عمودیت کی شرط سے، خط کی مساوات نکالنا۔
\end{enumerate}

\medskip
ان پانچ کو حل ڈھانپ کر متعلقہ یونٹ سے دہرائیے۔ پانچ مہارتیں جو بغیر دیکھے درست
آ جائیں، نو یونٹ سرسری پڑھ لینے سے کہیں زیادہ قیمتی ہیں۔

\vfill
\begin{center}
\begin{tcolorbox}[enhanced, colback=bmtint, colframe=bmblue, boxrule=0.5pt,
                  arc=3pt, width=0.92\linewidth,
                  left=12pt, right=12pt, top=10pt, bottom=10pt]
\small
\textbf{\color{bmblue}تفصیلی ایڈیشن۔}\ \
تمام نو یونٹوں کے مکمل حل، ہر سوال کی پڑتال، خاکے اور اضافی مشق کے لیے دیکھیے
\textenglish{AIOU\_1429\_Solved\_QuestionBank.pdf} (انگریزی،
\textenglish{84} صفحات)۔

\medskip
\textbf{\color{bmblue}ماخذ۔}\ \
\emph{کاروباری ریاضی}، کورس کوڈ \textenglish{5405/1429}، یونٹ
\textenglish{1}--\textenglish{9}، شعبۂ ریاضی، علامہ اقبال اوپن یونیورسٹی،
اسلام آباد، اشاعت \textenglish{2023}۔
\end{tcolorbox}
\end{center}
